% !TeX program = xelatex
% ============================================================================
%  第 4 课　群的四条规则　—— 教师版讲义
%  与 02-课件/第04课-群的四条规则/第04课-群的四条规则.tex 同步
%  记号：本课的新定义运算统一用 \oplus，避免在数学模式里排中文符号
% ============================================================================

\jhlesson{4}{群的四条规则}{时长 45 分钟\quad·\quad 教具：无\quad·\quad 本课不发单页}

\begin{mubiao}
  \begin{itemize}
    \item 能说出群的四条规则，并能在模 \(12\) 加法、有理数加法、新定义运算里各找到一个例子。
    \item 知道单位元不一定是 \(0\)（在 \(a\oplus b=a+b+1\) 里它是 \(-1\)）。
    \item 知道交换律不是群的必要条件，结合律也不是废话。
  \end{itemize}
\end{mubiao}

\noindent{\small 本课节奏：0—5 回顾｜5—12 群的四条规则（定义）｜12—25 四个例子逐条打勾｜25—34 结合律不是废话｜34—42 交换律不是必要的｜42—45 收尾。}

注意：「运算」的定义已在第 2 课讲过，本课不再重复，只做一句回顾，把省下来的时间给结合律反例。

\section*{一、回顾（0—5 分钟）}

把上节课的对照表重新写在黑板上（模 \(12\) 加法、有理数加法、有理数乘法，各自的单位元与逆元），
然后问一句：这三套运算看上去很像。\textbf{像在哪儿？}

今天就做一件事：把这个「像」写成一个定义。

\section*{二、群的四条规则（5—12 分钟）}

\begin{definition}[群]
  设 \(*\) 是集合 \(G\) 上的一个运算。如果下面四条都成立，
  就说 \((G,*)\) 是一个\textbf{群}：
  \begin{enumerate}
    \item \textbf{封闭}：任取 \(a,b\in G\)，\(a*b\) 仍然属于 \(G\)；
    \item \textbf{结合}：\((a*b)*c=a*(b*c)\)；
    \item \textbf{单位元}：存在 \(e\in G\)，对每个 \(a\) 都有 \(a*e=e*a=a\)；
    \item \textbf{逆元}：每个 \(a\in G\) 都有 \(b\in G\)，使 \(a*b=b*a=e\)。
  \end{enumerate}
\end{definition}

\begin{zhu}
  念完定义，立刻用手指敲一下黑板：「交换律不在里面。」
  \(a*b=b*a\) 不要求成立——这一点本课最后会专门拿出一个例子来说明。
\end{zhu}

\section*{三、四个例子，逐条打勾（12—25 分钟）}

板书这张表，一边填一边讲：

\begin{center}
\begin{tabular}{@{}lcccc@{}}
  \toprule
   & 封闭 & 结合 & 单位元 & 逆元 \\
  \midrule
  模 \(12\) 加法 & \(\checkmark\) & \(\checkmark\) & \(0\) & \(12-a\) \\
  有理数加法 & \(\checkmark\) & \(\checkmark\) & \(0\) & \(-a\) \\
  有理数乘法 & \(\checkmark\) & \(\checkmark\) & \(1\) & \(1/a\ (a\neq0)\) \\
  \(a\oplus b=a+b+1\) & \(\checkmark\) & \(\checkmark\) & \(-1\) & \(-a-2\) \\
  \bottomrule
\end{tabular}
\end{center}

第三行要停下来特别看一眼：\(a=0\) 的时候，\(1/a\) 是什么？——没有。

\begin{dingli}[有理数乘法不构成群]
  在有理数乘法里，\(0\) 没有逆元（要找 \(b\) 使 \(0\times b=1\)，而左边永远是 \(0\)）。
  少了一个逆元，四条规则就凑不齐，所以有理数乘法\textbf{不构成群}。
\end{dingli}

第四行是一个新定义的运算，当场带全班一条一条验：

\begin{example}[新定义运算 \(a\oplus b=a+b+1\)]
  \begin{itemize}
    \item 算两个试试：\(3\oplus5=3+5+1=9\)，\((-2)\oplus0=-2+0+1=-1\)。
    \item \textbf{封闭}：\(a\oplus b\) 是两个有理数相加再加 \(1\)，结果还是有理数。\(\checkmark\)
    \item \textbf{结合}：\((a\oplus b)\oplus c=(a+b+1)+c+1=a+b+c+2\)，
      而 \(a\oplus(b\oplus c)=a+(b+c+1)+1=a+b+c+2\)，两边一样。\(\checkmark\)
    \item \textbf{单位元}：解 \(a\oplus e=a\)，得 \(a+e+1=a\)，所以 \(e=-1\)。
      \textbf{注意：单位元是 \(-1\)，不是 \(0\)。}
    \item \textbf{逆元}：解 \(a\oplus b=-1\)，得 \(a+b+1=-1\)，所以 \(b=-a-2\)。
  \end{itemize}
\end{example}

\begin{caozuo}
  手指回答：在这个运算里，\(0\) 的逆元是 \(0\oplus b=-1\) 的解，即 \(-2\)；
  \(1\) 的逆元是 \(-3\)；\(-3\) 的逆元是 \(1\)。
\end{caozuo}

\section*{四、结合律不是废话（25—34 分钟）}

再换一个运算，这次故意让它有毛病：

\[
  a\oplus b = \frac{a+b}{2} .
\]

\begin{caozuo}
  当场算给全班看：
  \[
    (1\oplus3)\oplus5 = 2\oplus5 = \frac{2+5}{2}=3.5,
    \qquad
    1\oplus(3\oplus5) = 1\oplus4 = \frac{1+4}{2}=2.5 .
  \]
  两个结果不一样，让学生自己说出「结合律不成立」这四个字。
\end{caozuo}

\begin{zhu}
  这是整门课最值钱的一段，不要赶。让学生亲眼看到：
  结合律不是一句废话，它是可以被打破的。
\end{zhu}

\section*{五、交换律不是必要的（34—42 分钟）}

板书两个动作：\quad 动作甲 = 把数乘 \(2\)；\quad 动作乙 = 把数加 \(1\)。

取 \(x=1\)：

\begin{itemize}
  \item 先甲后乙：\(1\xrightarrow{\ \text{甲}\ }2\xrightarrow{\ \text{乙}\ }3\)；
  \item 先乙后甲：\(1\xrightarrow{\ \text{乙}\ }2\xrightarrow{\ \text{甲}\ }4\)。
\end{itemize}

\(3\neq4\)：先做哪个，结果不一样。

\begin{caozuo}
  接着问：那三个动作连着做呢？甲、乙、甲。
  先把前两个合成一个动作，再做甲；或者先做甲，再把后两个合成一个动作——
  两条路都是「甲、乙、甲」，结果一样。
\end{caozuo}

\begin{dingli}[复合的结合律]
  若干动作一个接一个地做，先把相邻的两个合成一个动作，不改变最后的结果。
\end{dingli}

\begin{zhu}
  结论落在这一句上：动作的复合天生满足结合律，却不一定满足交换律。
  这一组动作，连同它们复合出来的所有动作，构成一个\textbf{不交换的群}。
  交换律确实不在群的规则里。
\end{zhu}

\section*{六、收尾（42—45 分钟）}

板书三句话：群 \(=\) 一个集合 \(+\) 一个运算 \(+\) 四条规则（封闭、结合、单位元、逆元）；
单位元不一定是 \(0\)；交换律不在群的规则里。

预告：下节课去看三角形。它的动作也不交换。

\begin{xiaojie}
  \begin{itemize}
    \item 群 \(=\) 集合 \(+\) 运算 \(+\) 四条规则：封闭、结合、单位元、逆元。
    \item 单位元不一定是 \(0\)——在 \(a\oplus b=a+b+1\) 里它是 \(-1\)，逆元是 \(-a-2\)。
    \item 结合律不是废话（\((a+b)/2\) 就把它破坏了）；交换律不是必要的。
  \end{itemize}
\end{xiaojie}

\noindent\textbf{下节课}：三角形的对称——正三角形有 \(6\) 个对称动作，我们把它们写成两行式。
